% =======================================================
% 4. SPECIFICA DEI TEST
% =======================================================

\section{Specifica dei Test}
Questa sezione descrive la strategia di testing adottata dal gruppo Synergo Unit per garantire la conformità del sistema \textit{Second Brain} ai requisiti analizzati.

Lo stato del test è indicato come:
\begin{itemize}
    \item \textbf{I}: Implementato;
    \item \textbf{NI}: Non Implementato (pianificato per le fasi successive);
    \item \textbf{S}: Superato.
\end{itemize}

\subsection{Strategia di Test}
Per garantire la qualità del prodotto il gruppo adotta una strategia di test suddivisa per strato architetturale,
 utilizzando framework e tool specifici:
 \begin{itemize}
    \item Backend (Python): pytest.
    \item Frontend (JavaScript/Typescript):
        \begin{itemize}
            \item vitest per test unitari,
            \item eslint e prettier per analisi statica e formattazione.
        \end{itemize}
 \end{itemize}

I tool di analisi statica verranno eseguiti automaticamente in CI (GitHub Actions) per prevenire
 l'integrazione di codice non conforme.

\subsection{Test di unità (TU)}
Verificano il corretto funzionamento delle singole unità logiche del codice. Saranno implementati nella fase di Product Baseline\textsubscript{G} (PB).

%\begin{table}[H]
%    \centering
%    \renewcommand{\arraystretch}{1.3}
%    \begin{tabularx}{\textwidth}{|c|X|c|}
%        \hline
%        \rowcolor{primarycolor!10}
%        \textbf{ID} & \textbf{Descrizione} & \textbf{Stato} \\
%        \hline
%        TU1 & [Pianificato per fase PB] & NI \\
%        \hline
%    \end{tabularx}
%    \caption{Test di Unità}
%\end{table}

\subsection{Test di integrazione (TI)}
I test di integrazione hanno l'obiettivo di verificare la corretta comunicazione tra i diversi moduli del sistema.

Durante la fase RTB tali test non sono ancora previsti, poiché l'attività di sviluppo è limitata alla validazione della fattibilità tecnologica tramite il Proof of Concept. La pianificazione dettagliata dei test di integrazione verrà effettuata nella fase di Product Baseline\textsubscript{G} (PB), contestualmente alla definizione dell'architettura software definitiva. 

%
%\begin{table}[H]
%    \centering
%    \renewcommand{\arraystretch}{1.3}
%    \begin{tabularx}{\textwidth}{|c|X|c|}
%        \hline
%        \rowcolor{primarycolor!10}
%        \textbf{ID} & \textbf{Descrizione} & \textbf{Stato} \\
%        \hline
%        TI1 & [Pianificato per fase PB] & NI \\
%        \hline
%    \end{tabularx}
%    \caption{Test di Integrazione}
%\end{table}
%
\subsection{Test di sistema (TS)}
Verificano che il comportamento dell'intera applicazione sia coerente con quanto descritto nei Casi d'Uso\textsubscript{G} (UC).

\renewcommand{\arraystretch}{1.3}
\begin{xltabular}{\textwidth}{|c|X|c|c|}
    \caption{Mappatura Test di Sistema - Casi d'Uso} \label{tab:test-sistema} \\
    \hline
    \rowcolor{primarycolor!10}
    \textbf{ID} & \textbf{Descrizione} & \textbf{UC Rif.} & \textbf{Stato} \\
    \hline
    \endfirsthead
    
    \multicolumn{4}{c}%
    {{\bfseries \tablename\ \thetable{} -- Continua dalla pagina precedente}} \\
    \hline
    \rowcolor{primarycolor!10}
    \textbf{ID} & \textbf{Descrizione} & \textbf{UC Rif.} & \textbf{Stato} \\
    \hline
    \endhead

    \hline
    \multicolumn{4}{|r|}{{Continua nella pagina successiva...}} \\
    \hline
    \endfoot

    \hline
    \endlastfoot

    TS01 & Operazioni base: Scrittura testo, Copia, Taglia e Incolla. & UC 1-4 & NI \\ \hline
    TS02 & Formattazione\textsubscript{G} testo: Grassetto Inserimento, Applicazione e Rimozione. & UC 5-8 & NI \\ \hline
    TS03 & Formattazione testo: Corsivo Inserimento, Applicazione e Rimozione. & UC 9-12 & NI \\ \hline
    TS04 & Formattazione testo: Sottolineato Inserimento, Applicazione e Rimozione. & UC 13-16 & NI \\ \hline
    TS05 & Formattazione testo: Barrato Inserimento, Applicazione e Rimozione. & UC 17-20 & NI \\ \hline
    TS06 & Formattazione testo: Citazione Inserimento, Applicazione e Rimozione. & UC 21-24 & NI \\ \hline
    TS07 & Gestione Titoli: Inserimento, Applicazione e Rimozione. & UC 25-28 & NI \\ \hline
    TS08 & Gestione Link: Inserimento, Modifica e Rimozione link. & UC 29-31 & NI \\ \hline
    TS09 & Gestione Elenchi: Inserimento, Modifica, Disattivazione e Rimozione. & UC 32-39 & NI \\ \hline
    TS10 & Gestione Tabelle: Creazione, Modifica Righe e Colonne, eliminazione. & UC 40-46 & NI \\ \hline
    TS11 & Gestione Modifiche: Azioni di Undo\textsubscript{G} e Redo\textsubscript{G}. & UC 47, 48 & NI \\ \hline
    TS12 & Visualizzazione: solo Editor, solo Render\textsubscript{G} e Split View\textsubscript{G}. & UC 49, 50, 51 & NI \\ \hline

    TS13 & AI: Richiesta, Generazione, Visualizzazione, Accettazione, Rifiuto, Rigenerazione, Visualizzazione errore  riassunto. & UC 52-60 & NI \\ \hline
    TS14 & AI: Richiesta, Generazione, Visualizzazione, Accettazione, Rifiuto, Rigenerazione, Visualizzazione errore  traduzione. & UC 61-69 & NI \\ \hline
    TS15 & AI: Richiesta, Generazione, Visualizzazione, Accettazione, Rifiuto, Rigenerazione, Visualizzazione errore  riscrittura. & UC 70-78 & NI \\ \hline
    TS16 & AI: Apertura finestra, Richiesta, Generazione, Visualizzazione, Accettazione, Rifiuto, Rigenerazione, Visualizzazione errore  Distant Writing. & UC 79-88 & NI \\ \hline
    TS17 & AI: Richiesta, Generazione, Visualizzazione, Accettazione, Rifiuto, Rigenerazione, Visualizzazione errore  capello bianco. & UC 89-97 & NI \\ \hline
    TS18 & AI: Richiesta, Generazione, Visualizzazione, Accettazione, Rifiuto, Rigenerazione, Visualizzazione errore  capello rosso. & UC 98-106 & NI \\ \hline
    TS19 & AI: Richiesta, Generazione, Visualizzazione, Accettazione, Rifiuto, Rigenerazione, Visualizzazione errore  capello nero. & UC 107-115 & NI \\ \hline
    TS20 & AI: Richiesta, Generazione, Visualizzazione, Accettazione, Rifiuto, Rigenerazione, Visualizzazione errore  capello giallo. & UC 116-124 & NI \\ \hline
    TS21 & AI: Richiesta, Generazione, Visualizzazione, Accettazione, Rifiuto, Rigenerazione, Visualizzazione errore  capello verde. & UC 125-133 & NI \\ \hline
    TS22 & AI: Richiesta, Generazione, Visualizzazione, Accettazione, Rifiuto, Rigenerazione, Visualizzazione errore  capello blu. & UC 134-142 & NI \\ \hline

    TS23 & Persistenza: Creazione ed Eliminazione nota e salvataggio nota in formato Markdown. & UC 143-145 & NI \\ \hline
    TS24 & I/O: Download/Upload MD ed Esportazione (PDF/HTML\textsubscript{G}/JSON\textsubscript{G}). & UC 146-149 & NI \\ \hline
    TS25 & Gestione Utenti: Registrazione, Autenticazione e Logout Utente. & UC 150-152 & NI \\ \hline
    TS26 & Interazioni con DB: Visualizzazione, Salvataggio, Apertura Note da Remoto. & UC 153-155 & NI \\ \hline
    
\end{xltabular}

\subsection{Test di accettazione (TA)}
I test di accettazione hanno lo scopo di verificare la conformità del prodotto rispetto alle aspettative della proponente Zucchetti e ai requisiti definiti nell'Analisi dei Requisiti.

La loro definizione dettagliata verrà completata nella fase di Product Baseline\textsubscript{G} (PB), quando sarà disponibile una versione sufficientemente completa del prodotto da sottoporre a validazione.

%\begin{table}[H]
%    \centering
%    \renewcommand{\arraystretch}{1.3}
%    \begin{tabularx}{\textwidth}{|c|X|c|}
%        \hline
%        \rowcolor{primarycolor!10}
%        \textbf{ID} & \textbf{Descrizione} & \textbf{Stato} \\
%        \hline
%        TA01 & Gestione completa del ciclo di vita di una nota Markdown. & NI \\
%        \hline
%        TA02 & Corretta integrazione delle funzionalità di elaborazione AI. & NI \\
%        \hline
%        TA03 & Analisi multi-prospettiva funzionante (Metodo De Bono). & NI \\
%        \hline
%        TA04 & Importazione ed esportazione dati corretta e senza perdite. & NI \\
%        \hline
%    \end{tabularx}
%    \caption{Test di Accettazione}
%\end{table}